\documentclass[11pt,a4paper]{article}

% ======================================================================
%  Analyse mathématique du jeu « Catan » — version simplifiée
%  Mise au propre de etude.md
% ======================================================================

\usepackage[utf8]{inputenc}
\usepackage[T1]{fontenc}
\usepackage[french]{babel}
\usepackage{lmodern}
\usepackage{microtype}

\usepackage{amsmath,amssymb}
\usepackage{geometry}
\geometry{margin=2.4cm}
\usepackage{xcolor}
\usepackage{booktabs}
\usepackage{array}
\usepackage{enumitem}
\usepackage{graphicx}
\usepackage{caption}
\usepackage{tcolorbox}
\usepackage{hyperref}

\usepackage{tikz}
\usetikzlibrary{automata,positioning,arrows.meta,shapes.geometric,calc,backgrounds,fit}

% ----------------------------------------------------------------------
%  Couleurs des ressources
% ----------------------------------------------------------------------
\definecolor{resBois}{HTML}{2E7D32}   % bois
\definecolor{resBle}{HTML}{F2C037}    % blé
\definecolor{resMouton}{HTML}{8BC34A} % mouton
\definecolor{resMinerai}{HTML}{607D8B}% minerai
\definecolor{resArgile}{HTML}{C75B39} % argile
\definecolor{resDesert}{HTML}{E0D2B4} % désert
\definecolor{resOr}{HTML}{D4AF37}     % or

\definecolor{accent}{HTML}{1565C0}
\definecolor{accentdark}{HTML}{0D47A1}
\definecolor{softgray}{HTML}{F4F5F7}

\hypersetup{
  colorlinks=true,
  linkcolor=accentdark,
  urlcolor=accent,
  pdftitle={Analyse mathematique du jeu Catan},
  pdfauthor={Etude}
}

% ----------------------------------------------------------------------
%  Boîte « définition / remarque »
% ----------------------------------------------------------------------
\newtcolorbox{remarque}[1][]{
  colback=softgray, colframe=accent, boxrule=0.6pt, arc=2pt,
  left=8pt, right=8pt, top=6pt, bottom=6pt,
  fonttitle=\bfseries, title=#1
}

% Petite pastille de ressource
\newcommand{\pastille}[2]{\tikz[baseline=-0.6ex]\node[circle,fill=#1,text=white,
  font=\bfseries\footnotesize,inner sep=1.6pt,minimum size=1.4em]{#2};}

% Hexagone « pointy-top » de circonradius 0.95 centré en (#1,#2), rempli de #3
% (à utiliser dans un environnement tikzpicture ; nécessite la librairie calc)
\newcommand{\hexat}[3]{%
  \filldraw[draw=black!45, line width=0.6pt, line join=round, fill=#3]
    ($(#1,#2)+(90:0.95)$)  -- ($(#1,#2)+(30:0.95)$)  -- ($(#1,#2)+(330:0.95)$) --
    ($(#1,#2)+(270:0.95)$) -- ($(#1,#2)+(210:0.95)$) -- ($(#1,#2)+(150:0.95)$) -- cycle;}

% Boîte d'encodage (tuples) en police à chasse fixe
\newtcolorbox{encodage}{
  colback=softgray!60, colframe=black!25, boxrule=0.5pt, arc=2pt,
  left=8pt, right=8pt, top=5pt, bottom=5pt, fontupper=\ttfamily}

\title{\vspace{-1.0cm}\textbf{Analyse mathématique du jeu « Catan »}\\[2pt]
       \large Version simplifiée avec marché central}
\author{Étude}
\date{\today}

\begin{document}
\maketitle

\begin{abstract}
\noindent
Cette étude propose une analyse mathématique du jeu de société \emph{Catan} en vue
d'un apprentissage par renforcement (RL). Nous y introduisons une \textbf{version
simplifiée} du jeu, débarrassée des interactions directes entre joueurs, et dotée
d'un \textbf{marché central} ainsi que d'une nouvelle ressource, l'\emph{or}. Le
présent document décrit les règles de cette variante, le modèle de perception, la
machine d'état du jeu et la représentation du plateau sous forme de graphe.
\emph{Il s'agit du début de l'étude.}
\end{abstract}

\tableofcontents
\vspace{0.5cm}
\hrule
\vspace{0.5cm}

% ======================================================================
\section{Introduction}
% ======================================================================

\emph{Catan} est un jeu de société stratégique qui implique la collecte de
ressources, la construction de routes et de colonies, et l'interaction avec les
autres joueurs.

Pour rendre plus simple le système de commerce et de négociation, nous utilisons une
version simplifiée du jeu, possédant un \textbf{marché central} ainsi qu'une nouvelle
ressource d'\textbf{or}, qui pourra être utilisée pour acheter des ressources à un
taux variable.

Cette version simplifiée \textbf{retire les interactions directes entre les joueurs}
et se concentre sur la gestion des ressources et la planification stratégique, tout
en gardant l'essence du jeu original.

% ======================================================================
\section{Les cartes développement}
% ======================================================================

Dans un premier temps, les cartes développement seront grandement simplifiées. Le jeu
original comporte les cartes développement suivantes.

\subsection{Cartes du jeu original}

\begin{description}[leftmargin=1.4cm, style=nextline]
  \item[Chevalier (14 cartes)] Permettent de déplacer le voleur et de voler une
    ressource à un autre joueur. De plus, le joueur qui possède le plus de cartes
    Chevalier (après en avoir joué $3$) peut obtenir la carte \emph{« plus grande
    armée »}, qui vaut $2$ points de victoire.
  \item[Progrès (6 cartes)] Offrent des bonus tactiques en ressources ou en
    construction. Il en existe $3$ types, chacun en $2$ exemplaires :
    \begin{itemize}[nosep]
      \item \emph{Construction de routes} ($\times 2$) : placer gratuitement
            $2$ routes sur le plateau.
      \item \emph{Invention / Progrès scientifique} ($\times 2$) : prendre
            gratuitement $2$ cartes Ressource de son choix dans la banque.
      \item \emph{Monopole} ($\times 2$) : nommer une ressource ; tous les autres
            joueurs doivent donner toutes les cartes de cette ressource qu'ils
            possèdent.
    \end{itemize}
  \item[Point de Victoire (5 cartes)] Représentent des bâtiments importants de la
    colonie. Chaque carte vaut $1$ point de victoire. Contrairement aux autres,
    elles ne sont pas révélées à la pioche : elles restent \textbf{cachées} en main
    et ne sont dévoilées qu'à la fin de la partie, lorsque le joueur atteint
    $10$ points.
\end{description}

\subsection{Simplification retenue}

\begin{remarque}[Simplification des cartes développement]
Dans notre version simplifiée, nous \textbf{retirons les cartes Progrès}. Chaque
joueur dispose alors de points de victoire \emph{publics} et \emph{privés} :
\begin{itemize}[nosep]
  \item les points de victoire \textbf{publics} sont visibles par tous les joueurs ;
  \item les points de victoire \textbf{privés} restent cachés jusqu'à la fin du jeu.
\end{itemize}
Le fait de retirer les cartes Progrès simplifie la perception des joueurs pour le RL.
\end{remarque}

\subsection{Plus grande armée et route la plus longue}

Deux récompenses accordent chacune $2$ points de victoire \textbf{publics} et sont
\textbf{transférables} au cours de la partie.

\begin{description}[leftmargin=1.4cm, style=nextline]
  \item[Plus grande armée] Accordée au joueur ayant joué le plus de cartes Chevalier,
    dès lors qu'il en a joué au moins~$3$. Si un autre joueur le dépasse
    \emph{strictement} (en nombre de chevaliers joués), la récompense change de main.
  \item[Route la plus longue] Accordée au joueur possédant la plus longue chaîne
    continue de routes, d'au moins~$5$ segments. En cas d'égalité, le détenteur actuel
    conserve la récompense. Une chaîne ne peut pas \emph{traverser} une intersection
    occupée par une colonie ou une ville adverse : une colonie adverse peut donc couper
    une route, et faire perdre (ou réduire) la plus longue route.
\end{description}

\begin{remarque}[Calcul de la route la plus longue]
La longueur de la route d'un joueur est la longueur du plus long \emph{chemin sans
réemprunter une arête} (un \emph{trail}) dans le sous-graphe de ses routes, le passage
étant interdit à travers une intersection adverse. Une intersection adverse peut
\emph{terminer} une chaîne, mais pas la prolonger.
\end{remarque}

% ======================================================================
\section{Marché central}
% ======================================================================

Chaque joueur commence la partie avec $10$ d'or, et le marché central est initialisé
avec un prix de $10$ d'or pour chaque ressource. Le prix des ressources évolue en
fonction de l'\textbf{offre et de la demande} : les joueurs peuvent y acheter ou y
vendre des ressources.

\begin{remarque}[Loi d'évolution des prix]
À chaque ressource \textbf{achetée}, le prix de cette ressource augmente de $1$ d'or.
À chaque ressource \textbf{vendue}, le prix diminue de $1$ d'or. Le prix ne peut pas
descendre en dessous de $1$ d'or, mais peut augmenter sans limite.
\[
  p_{t+1} \;=\;
  \begin{cases}
    p_t + 1 & \text{après un achat,}\\[2pt]
    \max(1,\; p_t - 1) & \text{après une vente.}
  \end{cases}
\]
\end{remarque}

\paragraph{Exemple.} Si le prix du bois est de $10$ d'or et qu'un joueur achète
$2$ bois, le prix passe à $12$ d'or et le joueur a dépensé $10 + 11 = 21$ d'or. Si un
autre joueur vend ensuite $1$ bois, le prix retombe à $11$ d'or et le joueur reçoit
$11$ d'or pour la vente.

\paragraph{Affichage : prix d'achat et de vente.} Pour l'aspect visuel, chaque
ressource possède deux prix, vente et achat, séparés de $1$ d'or.

\begin{figure}[h]
\centering
\begin{tikzpicture}[font=\small]
  % --- Tableau "ardoise de marché" ---
  \node[draw=accentdark, fill=resBois, text=white, font=\bfseries,
        minimum width=3cm, minimum height=0.8cm, rounded corners=2pt] (h) {Bois};
  \node[draw=accentdark, fill=white, minimum width=3cm, minimum height=0.8cm,
        below=0pt of h] (a) {Achat \;\textbf{10}};
  \node[draw=accentdark, fill=white, minimum width=3cm, minimum height=0.8cm,
        below=0pt of a] (v) {Vente \;\textbf{9}};

  % --- Courbe d'évolution du prix ---
  \begin{scope}[shift={(6.5,-1.4)}]
    \draw[->] (0,0) -- (5.4,0) node[right,font=\footnotesize]{transactions};
    \draw[->] (0,0) -- (0,3.0) node[above,font=\footnotesize]{prix};
    \foreach \y in {1,2,3} \draw (-2pt,\y*0.7) -- (2pt,\y*0.7);
    \node[font=\scriptsize,left] at (0,0.7) {1};
    \node[font=\scriptsize,left] at (0,1.4) {10};
    \node[font=\scriptsize,left] at (0,2.1) {12};
    % achats : montée
    \draw[resArgile, very thick, -{Stealth}]
      (0.4,1.4) -- (1.4,1.4) -- (1.4,1.75) -- (2.4,1.75) -- (2.4,2.1);
    \node[resArgile,font=\scriptsize,align=center] at (1.6,2.55){achats\\($+1$)};
    % ventes : descente, bornée à 1
    \draw[resBois, very thick, -{Stealth}]
      (2.4,2.1) -- (3.4,2.1) -- (3.4,1.75) -- (4.6,1.75) -- (4.6,1.4);
    \node[resBois,font=\scriptsize,align=center] at (4.4,2.4){ventes\\($-1$)};
    \draw[dashed,gray] (0,0.7) -- (5.2,0.7);
    \node[gray,font=\scriptsize] at (4.3,0.55){plancher = 1};
  \end{scope}
\end{tikzpicture}
\caption{Affichage d'une ressource au marché (gauche) et dynamique du prix sous
l'effet des achats et des ventes (droite).}
\end{figure}

% ======================================================================
\section{Plateau et perception}
% ======================================================================

Nous distinguons deux types de perception dans le jeu : une perception
\textbf{fixe}, qui ne varie pas au cours de la partie, et une perception
\textbf{dynamique}, qui évolue en fonction des actions des joueurs.

\begin{table}[h]
\centering
\renewcommand{\arraystretch}{1.3}
\begin{tabular}{>{\raggedright\arraybackslash}p{6.6cm} >{\raggedright\arraybackslash}p{6.6cm}}
\toprule
\textbf{\textcolor{accentdark}{Perception fixe}} &
\textbf{\textcolor{accentdark}{Perception dynamique}} \\
\midrule
Disposition des tuiles et numéros associés à chaque tuile. &
Position des colonies, villes et routes des joueurs. \\
Disposition des ports et leur type. &
Position du voleur. \\
 & Ressources possédées par le joueur. \\
 & Nombre de ressources possédées par les autres joueurs. \\
 & Coût des ressources au marché central. \\
\bottomrule
\end{tabular}
\caption{Décomposition de la perception du joueur.}
\end{table}

Certaines informations ne seront pas ajoutées directement à la perception dynamique :
elles pourront plutôt être \textbf{déduites} par le joueur à partir de ses
observations. Nous \textbf{n'incluons pas} dans la perception dynamique :

\begin{itemize}[nosep]
  \item les ressources possédées par les autres joueurs ;
  \item les ressources disponibles dans la banque ;
  \item le nombre de cartes développement possédées par les autres joueurs ;
  \item le nombre de points de victoire des autres joueurs.
\end{itemize}

% ======================================================================
\section{Tour d'un joueur}
% ======================================================================

Les dés sont lancés, et les ressources sont distribuées en fonction du résultat et de
la position des colonies et villes des joueurs.

\begin{itemize}[nosep]
  \item Si le dé \textbf{n'a pas fait $7$}, le joueur joue son tour normalement.
  \item S'il a \textbf{fait $7$}, il déplace le voleur et vole une ressource à un
        autre joueur, puis joue son tour normalement.
\end{itemize}

\begin{remarque}[Défausse sur un $7$]
Lorsqu'un $7$ est obtenu, chaque joueur (y compris le joueur actif) possédant
\textbf{plus de $7$ cartes} ressource doit en défausser la moitié,
$\left\lfloor n/2 \right\rfloor$. Ces cartes sont \textbf{choisies aléatoirement} : le
joueur ne contrôle pas les cartes qu'il perd.
\end{remarque}

L'effet de déplacement du voleur est représenté par un état particulier de la machine
d'état, l'état \textbf{VOL}, qui peut aussi être déclenché lorsqu'un joueur joue une
carte Chevalier.

\paragraph{Actions possibles pendant le tour.}
\begin{itemize}[nosep]
  \item Acheter et vendre des ressources au marché central.
  \item Échanger des ressources avec les ports (ou $4$ pour $1$ avec la banque).
  \item Construire des routes, colonies et villes.
  \item Acheter des cartes développement.
  \item Jouer une carte Chevalier.
  \item Passer.
\end{itemize}

Cela permet de créer une routine \textbf{indépendante} des autres joueurs, tout en
gardant l'interaction \emph{indirecte} avec eux à travers le marché central et les
cartes Chevalier.

% ======================================================================
\section{Machine d'état du jeu}
% ======================================================================

Le jeu possède une machine d'état qui représente les différentes phases du jeu et les
transitions entre elles. Les états principaux sont :

\begin{itemize}[nosep]
  \item \textbf{PREP} — état de préparation : les joueurs placent leurs colonies et
        routes initiales.
  \item \textbf{DÉS} — lancer de dés : détermine quelles ressources sont produites et
        distribuées.
  \item \textbf{JOUEUR} — tour du joueur : le joueur actif effectue ses actions.
\end{itemize}

\subsection{Vue globale}

\begin{figure}[h]
\centering
\begin{tikzpicture}[
    >=Stealth, node distance=3.2cm, on grid,
    state/.style={circle, draw=accentdark, fill=softgray, very thick,
                  minimum size=1.5cm, font=\bfseries}]
  \node[state, initial, initial text=début] (prep) {PREP};
  \node[state, right=of prep] (des) {DÉS};
  \node[state, right=of des] (joueur) {JOUEUR};

  \path[->, thick]
    (prep)  edge node[above]{} (des)
    (des)   edge[bend left]  node[above]{distribution} (joueur)
    (joueur)edge[bend left]  node[below]{fin de tour} (des);
\end{tikzpicture}
\caption{Machine d'état globale : $\textbf{PREP} \to \textbf{DÉS}
\leftrightarrow \textbf{JOUEUR}$.}
\end{figure}

\subsection{Détail d'un tour de joueur}

\begin{figure}[h]
\centering
\begin{tikzpicture}[
    >=Stealth, node distance=2.6cm and 3.4cm, on grid,
    state/.style={rectangle, rounded corners=3pt, draw=accentdark, fill=softgray,
                  very thick, minimum width=2.6cm, minimum height=1cm,
                  align=center, font=\bfseries\small},
    vol/.style={state, fill=resArgile!25, draw=resArgile}]

  \node[state] (debut) {DÉBUT\\DE TOUR};
  \node[state, below=of debut] (jeux) {JEUX};
  \node[state, below=of jeux] (pass) {PASS};
  \node[vol, right=of jeux] (vol) {VOL};

  \path[->, thick]
    (debut) edge node[left, align=center]{}                       (jeux)
    (debut) edge[bend left=15] node[above, font=\scriptsize]{carte Voleur / dé = 7} (vol)
    (jeux)  edge[bend left=15] node[above, font=\scriptsize]{carte Chevalier} (vol)
    (vol)   edge[bend left=15] node[below, font=\scriptsize]{vol effectué} (jeux)
    (jeux)  edge node[left]{passer} (pass);
\end{tikzpicture}
\caption{Sous-machine du tour d'un joueur. L'état \textbf{VOL} est atteint
lorsqu'un $7$ est obtenu ou qu'une carte Chevalier est jouée.}
\end{figure}

% ======================================================================
\section{Représentation du plateau}
% ======================================================================

Le plateau est représenté par un \textbf{graphe} : les \emph{sommets} représentent
les intersections (où colonies et villes peuvent être construites) et les
\emph{arêtes} représentent les routes. Chaque tuile de ressource porte un numéro qui
détermine la production lorsque ce numéro est obtenu aux dés.

\subsection{Notation des ressources}

Pour simplifier, nous nommons les ressources comme suit.

\begin{table}[h]
\centering
\renewcommand{\arraystretch}{1.3}
\begin{tabular}{c l l}
\toprule
\textbf{Symbole} & \textbf{Ressource} & \\
\midrule
\texttt{B} & Bois            & \pastille{resBois}{B} \\
\texttt{W} & Blé (\emph{wheat})   & \pastille{resBle}{W} \\
\texttt{S} & Mouton (\emph{sheep})& \pastille{resMouton}{S} \\
\texttt{O} & Minerai (\emph{ore}) & \pastille{resMinerai}{O} \\
\texttt{C} & Argile (\emph{clay}) & \pastille{resArgile}{C} \\
\texttt{X} & Rien (désert ou mer) & \pastille{resDesert!80!black}{X} \\
\bottomrule
\end{tabular}
\caption{Symboles utilisés pour les ressources.}
\end{table}

\subsection{Illustration : tuiles, sommets et arêtes}

\begin{figure}[h]
\centering
\begin{tikzpicture}[scale=1.0,
    hexface/.style={draw=black!55, line width=0.7pt, line join=round},
    sommet/.style={circle, fill=accentdark, draw=white, line width=0.4pt,
                   inner sep=0pt, minimum size=4.6pt},
    route/.style={line width=3.2pt, resBois!70!black, line cap=round},
    token/.style={circle, fill=white, draw=black!45, line width=0.6pt,
                  minimum size=0.78cm, inner sep=0pt, font=\bfseries},
    res/.style={font=\bfseries\footnotesize}]

  % --- Sommets partagés (graphe G = (V,E)), tuiles « pointy-top » jointives ---
  %     pas horizontal = 1.126, vertical = 0.65 / 1.3 ; rangée du bas à y = -1.95
  \coordinate (a) at (-2.252, 0.65); \coordinate (b) at (-2.252,-0.65);
  \coordinate (c) at (-1.126, 1.30); \coordinate (d) at (-1.126,-1.30);
  \coordinate (e) at ( 0.000, 0.65); \coordinate (f) at ( 0.000,-0.65);
  \coordinate (g) at ( 1.126, 1.30); \coordinate (h) at ( 2.252, 0.65);
  \coordinate (i) at ( 2.252,-0.65); \coordinate (j) at ( 1.126,-1.30);
  \coordinate (k) at (-1.126,-2.60); \coordinate (l) at ( 0.000,-3.25);
  \coordinate (m) at ( 1.126,-2.60);

  % --- Faces des trois tuiles (jointives, sans recouvrement) ---
  \filldraw[hexface, fill=resBois!50]    (c)--(a)--(b)--(d)--(f)--(e)--cycle; % Bois
  \filldraw[hexface, fill=resBle!55]     (g)--(e)--(f)--(j)--(i)--(h)--cycle; % Blé
  \filldraw[hexface, fill=resMinerai!40] (f)--(d)--(k)--(l)--(m)--(j)--cycle; % Minerai

  % --- Étiquette de ressource + jeton numéro (centre de chaque tuile) ---
  \node[res, text=resBois!45!black]    at (-1.126, 0.46) {B};
  \node[token, text=red]               at (-1.126,-0.18) {8};
  \node[res, text=resBle!35!black]     at ( 1.126, 0.46) {W};
  \node[token, text=red]               at ( 1.126,-0.18) {6};
  \node[res, text=resMinerai!55!black] at ( 0.000,-1.49) {O};
  \node[token]                         at ( 0.000,-2.13) {11};

  % --- Sommets du graphe (toutes les intersections) ---
  \foreach \v in {a,b,c,d,e,f,g,h,i,j,k,l,m} \node[sommet] at (\v) {};

  % --- Une route construite (arête c--e) reliant une ville à une colonie ---
  \draw[route] (c) -- (e);

  % --- Ville (étoile) et colonie (carré) sur des sommets bien dégagés ---
  \node[star, star points=5, star point ratio=2.3, fill=resOr, draw=black!45,
        line width=0.4pt, minimum size=14pt, inner sep=0pt] at (c) {};
  \node[rectangle, fill=resArgile, draw=white, line width=0.5pt,
        minimum size=9pt, inner sep=0pt] at (e) {};

  % --- Annotations ---
  \node[font=\scriptsize, text=resOr!60!black, anchor=east]   at ($(c)+(-0.18,0.10)$) {ville};
  \node[font=\scriptsize, text=resArgile!85!black, anchor=west] at ($(e)+(0.18,0.18)$) {colonie};
  \node[font=\scriptsize, text=resBois!55!black, sloped, anchor=south]
        at ($(c)!0.5!(e)+(0,0.04)$) {route};

\end{tikzpicture}
\caption{Représentation du plateau : les tuiles hexagonales portent une ressource et
un numéro de production (en rouge pour les numéros fréquents $6$ et $8$). Les
\textbf{sommets} sont les intersections constructibles (colonie = carré, ville =
étoile) et les \textbf{arêtes} sont les routes.}
\end{figure}

\subsection{Coordonnées axiales}

Le plateau est une grille hexagonale ; nous utilisons donc un système de coordonnées
\textbf{axiales}. La position de chaque tuile est un couple $(q, r)$, où $q$ est la
coordonnée « horizontale » et $r$ la coordonnée « verticale ». La tuile $(0,0)$ est la
tuile \textbf{centrale} du plateau ; ses six voisines sont obtenues par les six
déplacements élémentaires.

\begin{table}[h]
\centering
\renewcommand{\arraystretch}{1.25}
\begin{tabular}{c l c l}
\toprule
$(q,r)$ & Direction & $(q,r)$ & Direction \\
\midrule
$(1,0)$  & droite              & $(-1,0)$ & gauche \\
$(1,-1)$ & haut à droite       & $(0,-1)$ & haut à gauche \\
$(0,1)$  & bas à droite        & $(-1,1)$ & bas à gauche \\
\bottomrule
\end{tabular}
\caption{Les six voisins de la tuile centrale $(0,0)$ en coordonnées axiales.}
\end{table}

\begin{figure}[h]
\centering
\begin{tikzpicture}[scale=1.0, font=\small]
  % --- 7 tuiles : centrale (0,0) + 6 voisines ---
  \hexat{0}{0}{resArgile!45}              % (0,0) centrale
  \hexat{1.645}{0}{softgray}              % (1,0)  droite
  \hexat{0.823}{1.425}{softgray}          % (1,-1) haut à droite
  \hexat{-0.823}{1.425}{softgray}         % (0,-1) haut à gauche
  \hexat{-1.645}{0}{softgray}             % (-1,0) gauche
  \hexat{-0.823}{-1.425}{softgray}        % (-1,1) bas à gauche
  \hexat{0.823}{-1.425}{softgray}         % (0,1)  bas à droite

  % --- étiquettes de coordonnées ---
  \node[font=\bfseries] at (0,0.34) {$(0,0)$};
  \node[token/.try, circle, fill=white, draw=black!45, minimum size=0.62cm,
        inner sep=0pt, font=\bfseries, text=red] at (0,-0.32) {6};
  \node[font=\scriptsize\bfseries, text=resArgile!70!black] at (0,0.62) {C};

  \node[font=\footnotesize] at ( 1.645, 0)     {$(1,0)$};
  \node[font=\footnotesize] at ( 0.823, 1.425) {$(1,\!-1)$};
  \node[font=\footnotesize] at (-0.823, 1.425) {$(0,\!-1)$};
  \node[font=\footnotesize] at (-1.645, 0)     {$(-1,0)$};
  \node[font=\footnotesize] at (-0.823,-1.425) {$(-1,1)$};
  \node[font=\footnotesize] at ( 0.823,-1.425) {$(0,1)$};
\end{tikzpicture}
\caption{Voisinage axial : la tuile centrale $(0,0)$ — ici une tuile d'argile
\texttt{C} portant le numéro $6$ — et ses six voisines.}
\end{figure}

\subsection{Représentation d'une tuile et d'un plateau}

Un plateau est une \textbf{liste de tuiles}, chaque tuile étant représentée par un
quadruplet $(q,\, r,\, \mathit{res},\, n)$ où $(q,r)$ sont les coordonnées axiales,
$\mathit{res}$ le type de ressource (cf. table~2) et $n$ le numéro associé à la tuile.

\begin{encodage}
(0, 0, C, 6)   \% la tuile centrale est une tuile d'argile portant le numéro 6
...
\end{encodage}

\subsection{Modèle de graphe : sommets et arêtes}

Pour représenter le placement des villes, colonies et routes, nous utilisons un graphe
$G = (V, E)$ où $V$ est l'ensemble des \textbf{sommets} (intersections) et $E$
l'ensemble des \textbf{arêtes} (routes reliant des intersections adjacentes). Chaque
sommet $v \in V$ peut être occupé par une colonie ou une ville appartenant à un joueur.

\paragraph{Encodage d'un sommet.} Chaque sommet $v \in V$ est décrit par :
\begin{itemize}[nosep]
  \item \textbf{trois couples} $(\mathit{r}, n)$ donnant les tuiles adjacentes
        (ressource et numéro). Un \textbf{port} apparaît sous la forme d'un couple
        $(\mathit{r}, -1)$ pour un port de ressource (échange $2{:}1$), ou
        $(\mathtt{X}, -1)$ pour un port générique $3{:}1$ ;
  \item un \textbf{indice de joueur} : le propriétaire de la colonie/ville sur ce
        sommet, ou $-1$ si le sommet est libre ;
  \item un \textbf{booléen} indiquant si le sommet porte une \emph{ville}
        (\texttt{true}) ou une \emph{colonie} (\texttt{false}).
\end{itemize}

\begin{encodage}
( (B,8), (W,4), (C,-1), -1, false )
\end{encodage}

\noindent
\emph{Lecture :} sommet \textbf{libre} ($-1$, \texttt{false}), adjacent à une tuile de
bois numéro $8$, à une tuile de blé numéro $4$, et à un \textbf{port d'argile} $2{:}1$
(couple $(C,-1)$).

\paragraph{Encodage d'une arête.} Chaque arête $e \in E$ est représentée par un triplet
$(v_1,\, v_2,\, p)$ où $v_1$ et $v_2$ sont les sommets reliés et $p$ l'indice du joueur
qui possède la route, ou $-1$ si l'arête est libre.

\begin{encodage}
( v1, v2, p )
\end{encodage}

\subsection{Conclusion : informations accessibles au joueur}

Le graphe $G$ remplace ainsi toute autre représentation du plateau. Pour conclure cette
partie, récapitulons les informations auxquelles un joueur a accès — c'est-à-dire son
\textbf{vecteur d'observation}, qui sera ensuite fourni au modèle d'apprentissage par
renforcement pour prendre ses décisions.

\begin{table}[h]
\centering
\renewcommand{\arraystretch}{1.35}
\begin{tabular}{>{\raggedright\arraybackslash}p{8.0cm} l}
\toprule
\textbf{Information observée} & \textbf{Représentation} \\
\midrule
Le graphe $G$ (plateau complet) & graphe \\
Ressources possédées par le joueur &
  \texttt{\{ressource: quantité\}} \\
Points de victoire \emph{publics} de chaque joueur &
  \texttt{[pv\_1, pv\_2, \dots]} \\
Cartes Point de Victoire possédées par le joueur &
  entier \\
Cartes Chevalier en main (privé) &
  entier \\
Nombre de ressources possédées par les autres joueurs &
  \texttt{[res\_1, res\_2, \dots]} \\
Prix des ressources au marché central &
  \texttt{\{ressource: prix\}} \\
\addlinespace[2pt]
\midrule
\multicolumn{2}{l}{\textit{Informations publiques pré-calculées}} \\
\midrule
Chevaliers \emph{joués} par chaque joueur &
  \texttt{[ch\_1, ch\_2, \dots]} \\
Longueur de la plus longue route de chaque joueur &
  \texttt{[lr\_1, lr\_2, \dots]} \\
Détenteur de la plus grande armée &
  entier ou $-1$ \\
Détenteur de la route la plus longue &
  entier ou $-1$ \\
\bottomrule
\end{tabular}
\caption{Vecteur d'observation accessible à un joueur.}
\end{table}

\noindent
On retrouve ici la distinction de la section~4 : le graphe $G$ porte la perception
fixe (tuiles, ports) et dynamique (colonies, villes, routes), tandis que les
ressources, les points de victoire publics et les prix du marché complètent l'état
observable. Les informations cachées (ressources \emph{détaillées}, cartes et points
\emph{privés} des adversaires) en sont volontairement absentes : elles ne pourront être
qu'\emph{estimées} par le joueur. Ce vecteur d'observation constitue l'entrée du modèle
d'apprentissage par renforcement.

\begin{remarque}[Pourquoi exposer la longueur des routes ?]
Les informations publiques pré-calculées sont en théorie \emph{déductibles} du graphe
$G$. Toutefois, déterminer la route la plus longue revient à résoudre un problème de
\textbf{plus long chemin} (combinatoire et global) : un modèle de régression — voire un
réseau de neurones sur graphe — ne saura pas le reconstituer de façon fiable à partir
de la seule structure de $G$. Comme ces informations sont \textbf{publiques}, les
pré-calculer et les exposer directement ne révèle aucune information cachée, tout en
épargnant au modèle un calcul qu'il ne maîtrise pas. La longueur des routes, le nombre
de chevaliers joués et les détenteurs des deux récompenses sont donc fournis
explicitement.
\end{remarque}

\vfill
\begin{center}\itshape\small
Fin de cette première partie de l'étude.
\end{center}

\end{document}
